\section*{Chapter \ref{ch:g6:solutions-and-solubility} --- Solutions and Solubility}

\begin{solution}{exo:g6:solutions-and-solubility:1}
\omterm{def:g6:solutions-and-solubility:solute}{Solute}: sugar. \omterm{def:g6:solutions-and-solubility:solute}{Solvent}: water. Yes, the \omterm{def:g6:solutions-and-solubility:solute}{solvent} is water: it is an
\omterm{def:g6:solutions-and-solubility:solute}{aqueous solution}.
\end{solution}

\begin{solution}{exo:g6:solutions-and-solubility:2}
A solution that cannot \omterm{def:g2:mixing-and-dissolving:dissolve}{dissolve} any more of its \omterm{def:g6:solutions-and-solubility:solute}{solute}. \omterm{def:g6:solutions-and-solubility:solute}{Solute} added to
it stays undissolved at the bottom.
\end{solution}

\begin{solution}{exo:g6:solutions-and-solubility:3}
$250 + 20 = \qty{270}{g}$.
\end{solution}

\begin{solution}{exo:g6:solutions-and-solubility:4}
Ethanol: \omterm{def:g6:solutions-and-solubility:miscible}{miscible}. Cooking oil: \omterm{def:g6:solutions-and-solubility:miscible}{immiscible}. Vinegar: \omterm{def:g6:solutions-and-solubility:miscible}{miscible} (it is
mostly water).
\end{solution}

\begin{solution}{exo:g6:solutions-and-solubility:5}
Propanone (acetone). GHS02: highly flammable; GHS07: irritates the eyes,
the vapour can cause drowsiness or dizziness.
\end{solution}

\begin{solution}{exo:g6:solutions-and-solubility:6}
$360 \times \frac{500}{1000} = \qty{180}{g}$.
\end{solution}

\begin{solution}{exo:g6:solutions-and-solubility:7}
\qty{200}{g} of water \omterm{def:g2:mixing-and-dissolving:dissolve}{dissolve} at most $360 \times \frac{200}{1000} =
\qty{72}{g}$. No: $80 - 72 = \qty{8}{g}$ stay at the bottom.
\end{solution}

\begin{solution}{exo:g6:solutions-and-solubility:8}
The solution weighs $225 + 25 = \qty{250}{g}$; $\frac{25}{250} =
\qty{10}{\percent}$.
\end{solution}

\begin{solution}{exo:g6:solutions-and-solubility:9}
In the first beaker, where all the salt has dissolved and the solution is
not yet saturated.
\end{solution}

\begin{solution}{exo:g6:solutions-and-solubility:10}
The carbon dioxide was dissolved under pressure in the closed bottle.
Opening it lowers the pressure: the water can hold less gas, and the
extra gas comes out as bubbles.
\end{solution}

\begin{solution}{exo:g6:solutions-and-solubility:11}
$2000 \div 0.04 = 50\,000$ times more.
\end{solution}

\begin{solution}{exo:g6:solutions-and-solubility:12}
A gas \omterm{def:g2:mixing-and-dissolving:dissolve}{dissolves} less in warm water: the warm pond holds less dissolved
oxygen, and the fish lack oxygen.
\end{solution}

\begin{solution}{pb:g6:solutions-and-solubility:1}
\textbf{1.} \omterm{def:g6:solutions-and-solubility:solute}{Solute}: salt. \omterm{def:g6:solutions-and-solubility:solute}{Solvent}: water.

\textbf{2.} The water holds as much salt as it can: no more salt can
\omterm{def:g2:mixing-and-dissolving:dissolve}{dissolve} in it.

\textbf{3.} The undissolved salt at the bottom proves that the brine is
saturated: it holds as much salt as possible.

\textbf{4.} $360 \times 20 = \qty{7200}{g} = \qty{7.2}{kg}$.

\textbf{5.} $20 + 7.2 = \qty{27.2}{kg}$.

\textbf{6.} $\frac{7.2}{27.2} \approx 0.26$: about \qty{26}{\percent}.

\textbf{7.} No: the water could \omterm{def:g2:mixing-and-dissolving:dissolve}{dissolve} \qty{7.2}{kg}. It could still
\omterm{def:g2:mixing-and-dissolving:dissolve}{dissolve} $7.2 - 5 = \qty{2.2}{kg}$.

\textbf{8.} $20 - 5 = \qty{15}{kg}$ of water (\qty{15}{L}).

\textbf{9.} $360 \times 15 = \qty{5400}{g} = \qty{5.4}{kg}$.

\textbf{10.} It comes out of solution as solid crystals: it
crystallises and sinks to the bottom.

\textbf{11.} New white crystals appear and pile up at the bottom of the
tank.

\textbf{12.} $7.2 - 5.4 = \qty{1.8}{kg}$ of salt have crystallised.
\end{solution}
